\chapter{Hitting times and the reciprocal xi function}\label{chap:hitting}\label{ch:ht}\chaptermark{Hitting times and reciprocal xi}

The reciprocal of the normalised $\xi$ function is the Laplace transform of a probability law, the law of the clock $T$ of Section~\ref{ch:clock}. Under the Riemann Hypothesis that law is the hitting time of a Krein string, the explosion time of a pure birth chain, and its density is Kent's eigenfunction expansion over the zeros. This chapter asks whether the argument can be run backwards: whether RH can be obtained from the limiting behaviour of hitting times. It shows where the content of such an argument has to lie. It tests the families of approximating laws that can be built without assuming anything about the zeros (Jensen laws, Gram smoothed zeros, truncated Euler products, boundaries of Brownian motion), proves a universal stable limit for the prime side hitting times and a rigidity theorem for factorisations of the clock law, and studies the density of the clock directly: its residue expansion, its log concavity, and the failure of hyperbolic complete monotonicity. No step proves RH.

The chapter develops the source chapter on hitting times and two notes, one on a residue expansion of the reciprocal $\xi$ function and one on the Euler product as a L\'evy measure, all by Nick Polson \cite{P26ht,PolsonResidue26,PolsonEuler26}. Its numbers were recomputed for this book by scripts written from the mathematics alone (the scripts of the source were not used); every number printed below is one of those recomputations unless it is marked unchecked, and where the recomputation differed from the source the recomputed value is printed. Results that no independent computation reached are said to be unchecked. The notation is the book's: the entire function is $G(u)=\xi(\tfrac12+\sqrt u)/\xi(\tfrac12)$ of Section~\ref{ch:clock}, the clock is $T$ with law $\mu$ and Laplace transform $1/G$, the canonical function is the heat trace $W$ (Chapter~\ref{chap:heat}), and the kernel of de Bruijn and Newman is $\Phi_{\mathrm{dBN}}$ of Section~\ref{ch:dbn}, not the kernel $\Phi$ of \eqref{eq:ch2:Phi}.

Several parts of the sources are already in the book and are cited, not restated: the self decomposability of $\mu$ and the equivalence of the GGC property with RH (Theorem~\ref{thm:ch12:idsd} and Proposition~\ref{prop:rigid}), the spectral expansion of the density and the Kent approximants (Section~\ref{ch:clocklaw}), the criterion by P\'olya frequency functions (Proposition~\ref{tx2:prop:pf}), the Jensen ladder (Proposition~\ref{pol:prop:lp}), the closure of the GGC class under weak limits (Corollary~\ref{cor:weak:closure}), the sine squared lemma (Lemma~\ref{rz:lem:sin2}) and the monotonicity ladder of the heat trace (Section~\ref{ch:heat}).

Status words follow the rest of the book. A statement is proved when a proof or a pointer is given, computed when a script of the book produced it (code/README.md lists the scripts \texttt{ht\_*.py} of this chapter), conditional, heuristic or open as stated. The results are these.
\begin{enumerate}
\item With no hypothesis on the zeros the density of $\mu$ is the limit of grouped residue sums, and its tail is $w_1e^{-\gamma_1^2y}+w_2e^{-\gamma_2^2y}+O(e^{-c_3y})$ (Theorem~\ref{ht:thm:grouped}, proved, and Proposition~\ref{ht:prop:tail}, proved). Absolute convergence of the residue series needs a lower bound for $|\xi'(\rho)|$ and is conditional (Theorem~\ref{ht:thm:res}). Under RH the partial fraction series of $1/G$ does not converge absolutely, because the weights $|w_\gamma|$ grow like $e^{\pi\gamma/4}$ (Proposition~\ref{ht:prop:krein}, proved).
\item Truncating the Hadamard product gives GGC laws exactly while the zeros used lie on the line, and $\mu_N$ is a GGC for every $N$ up to the zeros below $3\cdot10^{12}$ (Proposition~\ref{ht:prop:trunc}, proved given the verification of RH to that height \cite{PT21}). The density is not HCM without RH, which is proved (HCM implies GGC, and GGC implies RH); the sum of the first $120$ residue terms vanishes at a point $z_*$ off the positive axis (computed), which would prove that the density is not HCM if the omitted terms are controlled by an explicit remainder bound on a sector containing $z_*$ (Theorem~\ref{ht:thm:nohcm}), and non HCM is proved for every finite sum of two or more exponentials (same theorem).
\item The density of $\mu$ is log concave in the tail (Theorem~\ref{ht:thm:lctail}, proved), on $[0.008,\infty)$ (Theorem~\ref{ht:thm:lc}, computed, with one numerical input not enclosed) and near the origin conditionally on a saddle point approximation (Proposition~\ref{ht:prop:lc0}, conditional).
\item The Jensen polynomials of $G$ give GGC laws $\mu_d$ for every $d\le2\cdot10^{17}$ (Proposition~\ref{ht:prop:jensen}, proved given a theorem of Chasse that was not read), converging to $\mu$ at the rate $2.8/d$ in distribution function (computed, on a grid of $[0.007,0.08]$). The resolution law $d^*=\gamma^2/(2t_c)$ for an off line pair is heuristic, and it reproduces the directly tracked Davenport and Heilbronn value to $0.1\%$ (computed) (Section~\ref{ht:sec:jensen}).
\item Every truncation of $\mu$ with real knots (zeros on the line, or under RH) is a B spline transform and a linear functional of the Gumbel extremal process (Proposition~\ref{ht:prop:bspline} and Theorem~\ref{ht:thm:extremal}, proved). The smoothed zero laws are within $1.3\%$ of $\mu$ in distribution function with no zero and within $2.5\cdot10^{-5}$ with ten zeros (computed). The P\'olya density is not $\mathrm{PF}_5$ (computed).
\item The truncated Euler hitting times satisfy $T_{\E\log N_P}\overset d=Y_P+T_{\log N_P}$ with $Y_P$ a GGC, and after centring and scaling they converge to a spectrally positive stable law of index $\frac32$ that does not depend on the zeros (Proposition~\ref{ht:thm:trunc} and Theorem~\ref{ht:thm:32}, proved; the rate is computed).
\item Factorisations $\mu=\mu_1*\mu_2$ into two GGC laws are rigid (Theorem~\ref{ht:thm:rigid}, proved). What this and the stable limit say about a hitting time proof built from the primes is a heuristic and is stated as one (Remark~\ref{ht:cor:shape}).
\item Hitting times of diffusions started at an entrance or reflecting end are GGC (Kent, proved); the identification of the limiting diffusion is open (Section~\ref{ht:sec:diffusion}).
\end{enumerate}
The later sections treat curved boundaries and the inverse first passage problem (Section~\ref{ht:sec:rs}), Yor's probabilistic Riemann function (Section~\ref{ht:sec:yor}), the Davenport and Heilbronn test (Section~\ref{ht:sec:dh}), the frontier of what the reciprocal $\xi$ law proves (Section~\ref{ht:sec:frontier}), the ladder of Taylor coefficients, what a hitting time proof would need, and open problems.
